\quellseite{126}
\part{Part IV. Theoretical and Practical}

\quellseite{128}
\chapter[Chapter XI. The Rite of Confirmation]{Chapter XI. The Catechist and the Rite of Confirmation}

CHAPTER XI.

THE CATECHIST AND THE RITE OF CONFIRMATION.

The rite of Confirmation closes the Catechist’s work of instruction. Not that he is now done with his class. He will never be done with a single one of them unless that one removes into another parish, or dies, or he takes charge of another parish. His catechising and confirming of the Catechumen make him, in a peculiar sense, his beichtkind, his spiritual child, to whom he is to stand in a specially interesting sense in the relation of a spiritual guide and father. But the Confirmation closes the regular and stated lessons and recitations of that class.

What is the meaning of Confirmation? What does it mean for the Catechist? What for the catechumen? \randnotiz{Confirmation—meaning}

We consider it here not in the broad sense as including, 1, Instruction; 2, Confession, and 3, Laying on of hands. We confine ourselves here to the narrow sense and consider only the act of public confession followed at once by the imposition of hands and the prayers of pastor and people.

In the Roman Catholic and Greek Catholic churches it is a sacrament. The ceremony is accompanied with the anointing with oil, and the confirming
\quellseite{129}
bishop says: “I seal thee with the sign of the cross and confirm thee with the oil of salvation, in the name of the Father, and of the Son, and of the Holy Ghost.” The effect of the crossing, anointing and imposition of hands is the giving of the Holy Ghost in a sense and measure in which He was not given at baptism. The sacramental character of Confirmation is higher than that of baptism. An ordinary priest may baptize, but only a bishop can confirm.

The Lutheran Reformers unanimously rejected \randnotiz{Not a Sacrament} the sacramental character of Confirmation. Luther, in his “Babylonian Captivity,” expresses himself as willing to retain it as a useful and edifying ceremony of the Church, if only it be freed from the mummeries and superstitions connected with it by the Romanists. He will not for a moment admit that it is a Sacrament. Both in the Augsburg Confession and in the Apology it is treated as an ecclesiastical rite, devoid of sacramental character. The Interims of 1548, which were Romanizing in their contents and tendency, enjoined it. But it was not generally practiced. Instruction in the Catechism was general and was faithfully practiced. An examination in the knowledge and understanding of the Catechism was usual before the admission to the first Communion. The Saxon General Articles of 1557 call this the true Confirmation.
\quellseite{130}
Bugenhagen introduced our present form of Confirmation into Pomerania. But it did not meet with universal favor. After Bugenhagen’s time it again fell into disuse, although we find it commended by such theologians as Gerhardt, Hunnius, Hutter, and others.

To Spener and his followers belongs the credit of making the act of Confirmation popular and general in the Lutheran Church. As we have already seen, Spener wanted it preceded \randnotiz{Spener and Confirmation} by a thorough catechisation and an impression on the heart and will by the class instruction. In this he was certainly correct. It is this failure to emphasize the influence of the acquired truth on heart and will that has so often and so sadly brought Confirmation into disrepute. True, some of the later Pietists were not as sound as Spener. They tried to bring in a false and forced emotionalism. Instead of so applying the catechetical truth that it should move the heart and will, they introduced unhealthy, Methodistic pressure and often tried to make Confirmation mean a sudden, conscious conversion. But we find nothing of this in Spener’s Explanation of Luther’s Catechism. The unsound extremes of a later age are not true Pietism, but its unnatural perversion.

An unsound Pietism was also unclear on the relation of Confirmation to Baptism. It was often represented as adding something to Baptism. \randnotiz{Unsound Pietism} It
\quellseite{131}
was forgotten that baptism is an act of God. As such it is complete. A human ceremony cannot add to or complete it. Again it was represented as a renewing of the baptismal covenant. Now the covenant includes what God has done, and God’s part cannot be renewed by man. But a covenant demands two parties. One gives, another receives; one proposes, another agrees; one conditions, another accepts. And so God and the baptized child are brought into covenant relation with one another. And the child neither renews nor ratifies the whole covenant.

But in the making of that covenant the child could not consciously take part. God could give, but the child could not consciously and verbally receive. Parents or sponsors stood in the child’s place and promised and covenanted in the child’s stead. Now, after due instruction \randnotiz{Confirmation—how related to baptism} as to the meaning of the baptismal act and covenant, the child answers in his own name, and in this sense takes upon himself and ratifies his side of the covenant. The whole preceding instruction was calculated to nourish and develop the germ-life implanted in baptism. The mind has been enlightened, the heart moved and drawn, and the will persuaded. The child is thus equipped consciously and intelligently to assume and ratify his side of the baptismal covenant. In this sense Confirmation looks back to and is related to baptism.
\quellseite{132}
But an unsound Pietism was not alone in misunderstanding and misrepresenting baptism. Long before Spener, that unstable, uncertain, mediating character, Bucer, for whom the Augsburg Confession was not good enough, who wrote the Tetropolitan Confession as a substitute, \randnotiz{Bucer} who was now on the Lutheran side and again on the side of the Reformed, gave a Romanizing interpretation to Confirmation. He introduced in Hesse the Formula in which the minister, as he lays his hand upon the child says: “Receive the Holy Ghost, safeguard and shelter against all malice, strength and help toward all good, from the gracious hand of God the Father, the Son and the Holy Ghost. Amen.”


Andreas Hyperius, according to Achelis,\footnote{“Practische Theologie,” p. 38.} teaches that while in baptism \randnotiz{Andreas Hyperius} only the forgiveness of sins is given, in Confirmation, by the laying on of hands, the Holy Ghost is given.


Vilmar, a noble, consecrated, confessional Lutheran of last century, who harbored sacerdotal views of the ministry, also taught that in Confirmation \randnotiz{Vilmar} the Holy Ghost is given by the laying on of hands. This Romanizing idea is closely connected with his hierarchical idea of ordination and the ministry.\footnote{See “The Lutheran Pastor,” p. 109 f.}


Others who do not share these views of Vilmar on the ministry yet seem to place a sort of a magical
\quellseite{133}
effect in the laying on of hands. They seem to believe and teach that on whomsoever the pastor’s hands are laid he is by virtue of this magic touch, without further consideration, fit and qualified to receive the Lord’s Supper.


All this we believe to be thoroughly unscriptural, \randnotiz{Laying on of hands} Romish and ruinous to souls. There is not the least warrant for the theory that grace can be conferred by the hands of pastor, priest or bishop. The miraculously endowed Apostles could confer charismatic gifts by the laying on of hands. Their uninspired and unendowed followers could not. The laying on of the confirming hands can at best indicate that the prayers, the promises and the blessings accompanying the act are surely meant for him whom the hand touches. In that sense the hands are aids for the subjective reception of the grace prayed for, promised and conferred by the Word. It is here as everywhere the Word that offers the grace and actually confers it where the recipient does not resist, but submissively and believingly allows that spirit-bearing Word to enter. This is the only evangelically-admissible view of the meaning of the laying on of hands, and this is enough.

As we have elsewhere written:\footnote{“Way of Salvation in the Lutheran Church,” p. 93 ff.} “This simple and appropriate ceremony we call Confirmation. We claim for it no magical powers. It is not a
\quellseite{134}
Sacrament. It adds nothing to the Sacrament of baptism, for that is complete in itself. The confirming, strengthening and establishing \randnotiz{Confirmation—what it is} of the catechumen in grace is effected primarily alone through Christ’s own means of grace, the Word and the Sacraments. The Word has been applied to mind and heart all along from tenderest childhood. It is now brought home in the review and admonition of the pastor, amid specially solemn surroundings. The previous administration of baptism and the perpetual efficacy of that Sacrament are now vividly recalled and impressed. And this unusually impressive application of the power of Word and Sacrament confirms and strengthens the divine life in the catechumen. Thus the means of grace do the confirming, or rather the Holy Spirit through these means. Instrumentally the pastor may be said to confirm, since he, as Christ’s ambassador or agent, applies His means of grace.


“In still another, though inferior sense, the catechumen confirms. He receives the offered means of grace, assents to their truth and efficacy, obtains divine virtue and strength through them, and with this imparted strength lays hold on Christ, draws nearer to Him, is united to Him as the branch to the vine, and thus confirms and establishes the covenant and bond that unites him to his Saviour.

“We do not claim for this rite of Confirmation
\quellseite{135}
a ‘thus saith the Lord.’ We do not claim that it possesses sacramental efficacy or that it is essential to salvation. . . . We believe it to be in perfect harmony with the whole tenor and spirit of the gospel. . . . We can find it in all its essential features in the pure age of the Church immediately succeeding the Apostles. In some form or other it has been practiced in the Church ever since.”

Like many other good institutions this one also has been abused. In German and Scandinavian countries it has shared in the inevitable evils of the State Church. Where the State made it compulsory it often degenerated into a mere formality that must be gone through with; the sooner and the easier the better.\footnote{Read carefully Chapter XIII., “The Lutheran Pastor.”}

“Being necessary to citizenship, \randnotiz{Abuse of Confirmation} it became also a badge of respectability and a mark of one’s standing in the community. To neglect it was not only a civil misdemeanor, but also a social disgrace. . . . Confirmation is regarded as a kind of a graduation from religious instruction. The confirmant is promised a new suit, a new dress, a watch, or both. After Confirmation there is a feast for those confirmed, and it is not unheard of that they have a social dance. Where such views and customs prevail, the children naturally want to be confirmed that they may be through with the drudgery of preparation and get
\quellseite{136}
the benefits that accompany and follow the ceremony.”

Unfortunately many of these abuses connected with Confirmation have been transplanted from the State churches of Europe to our land.

Here, too, in many of the foreign-speaking churches, we find the graduation idea; the being done with instruction and learning idea; the tendency to give up Sunday school and church attendance after Confirmation; the keeping the mind and heart of the catechumen on the new clothes, the dress parade; the laying of the chief emphasis on intellectual acquirements and the neglect of the heart. To some of these things we shall return when we consider what constitutes fitness for Confirmation.

The above abuses are by no means constituent parts of Confirmation. They are worldly accessories, brought into connection with this blessed rite in worldly periods, by worldly congregations and pastors. They cause our earnest English Lutheran pastors much perplexity and heart sorrow. They have much to do with the serious losses of the confirmed ones. We know of foreign-speaking congregations where those confirmed were not considered members of the congregation. To become such they were required to make special application at some subsequent time.

When freed from these abuses, so foreign to the
\quellseite{137}
true idea of Confirmation, \randnotiz{Confirmation a blessed ceremony} it is one of the most beautiful, blessed and impressive ceremonies of the Church. There can be no more befitting, touching and solemn ceremony of admitting the children of the Church into full communicant membership than Confirmation, when rightly conducted. Who has not witnessed this never-to-be-forgotten ceremony? And what could be more interesting, significant and soul-moving than to see a band of young, well-instructed and consecrated disciples encircling the altar of the Church, confessing their faith and bowing their knees to their Saviour amid the prayers and benedictions of the congregation? In concluding this chapter we quote again from “The Lutheran Pastor”:\footnote{P. 318 f.}

“Let the pastor prepare himself thoroughly for the solemnities of Confirmation day. It ought to be a day never to be forgotten. A festive solemnity and a thankful seriousness should pervade the whole service. Everything should be clearly understood by the catechumens. The girls should come forward and kneel without hats. If any are to be baptized they should be placed at the end of the row. All should be in the places reserved for them, in good time. The hymns, lessons and prayers should be suited to the occasion. The sermon should breathe the spirit of loving solicitude for the young
\quellseite{138}
disciples. It should contain most earnest and pointed warnings for those in the audience who have forgotten their Confirmation vows. Parents and friends should be admonished of their duty to help the confirmed to remain true. All the members should be exhorted to welcome into their fellowship, to pray for and help these new communicants. The formula should be read slowly, distinctly and impressively. Before the free prayer for these particular youths it would be well to ask the whole congregation to pray also. Let there be no haste. It is a high day in the lives of these dear ones. It ought to be a high day for the congregation. Ah, these Confirmation days! How we recall their prayers and benedictions and tears! Shall they ever be forgotten? Should not the memory of them be like the ringing of heaven’s chimes to recall the wanderer? Let the pastor make the most of them.”

\quellseite{139}
\chapter{Chapter XII. The Catechist’s Qualifications}

CHAPTER XII.

THE CATECHIST’S QUALIFICATIONS.

Model Catechists, like model pastors, are rare. Men may be good preachers and poor Catechists. They may be successful in pastoral work and fail in their catechetical functions. Qualifications are needed. Some are natural, others can be acquired. All can be improved. We can divide the qualifications into the personal and the professional.

What, then, must the good Catechist be? \randnotiz{A good Catechist} And what must he know? He must be, first of all, a devoted Christian. It is not enough that he be an ordinary Christian. As more is expected of him than of the average everyday Christian, so more is required of him in his personal, spiritual character.

One thing is very significant and should be matter of the deepest concern to him who would excel as a Catechist. And right here let us again emphasize that every minister of Christ should have a burning zeal, a sanctified ambition, an ambition purged from self-seeking, to excel in His service, to be the very best preacher, seelsorger and Catechist that it is possible for him to be. He should covet earnestly the best gifts. A minister who is satisfied, if he can carry along his work in the easiest
\quellseite{140}
possible way, satisfied if he can hold his place and have his bread, will have a heavy account to render. The one significant and serious thing for his prayerful consideration is the fact that before Jesus reinstated the fallen Peter, before He commissioned him to feed His lambs, He asked him to answer the searching question, “Lovest thou me more than these?” The Master repeated the question. He probed deeply, \randnotiz{First qualification—love to Jesus} and not until Peter, humbled and grieved, looked up into His face and out of the innermost depths of his warm heart repeated, “Lord, Thou knowest all things: Thou knowest that I love Thee,” was Jesus satisfied to have him officially feed His lambs, His sheeplings, His sheep. Love to Jesus; personal love; experimental love; love out of a penitent, believing, longing heart; love like Peter’s, felt in the heart, confessed with the mouth, ready to serve, ready to suffer: this is the first and great qualification. Without love all other gifts and graces are sounding brass, tinkling cymbal, nothing. Out of a heart, broken by penitence, healed by faith, burning with love, with upturned face and soulful voice he must be able to sing:


“I’ll go where you want me to go, dear Lord,
I’ll be what you want me to be.”

Another qualification of the good Catechist is that he must be a lover of children. There are men who
\quellseite{141}
have an aversion to children. \randnotiz{Second—a lover of children.} They are uncomfortable when children are near. They have only annoyed, unkind, if not scowling looks for them. It is a question whether men who have and harbor such a nature ought ever to become fathers. They are certainly out of place in the office one of whose chief functions it is to feed Christ’s lambs.

The good Catechist finds pleasure in children. He delights to gather them around him. Their liveliness, their curiosity, their prattle, their very restlessness and mischief are sources of unfailing interest to him. He finds joy in studying them, in winning them, in entertaining them. He has patience with their childishness. He directs and helps their impulsiveness. He enters into their interests. He becomes a child among children. He has a warm place in his heart for the boy. He does not censure the boy for brimming over with vitality and energy. He knows that this is a part of the healthy boy’s nature. And when that boy’s restlessness and mischief and explosive energy are trying and vexing in themselves the good Catechist still loves him. Instead of shunning him, he cultivates him and studies how to win him. He loves children. He loves that unique, vivacious, mischief-loving and trying species—the American boy. He studies that boy, he chums with him, he wins him, he directs that ebullient energy. He makes a man of him. Blessed
\quellseite{142}
is the Catechist who knows and loves and wins and holds the children, especially the boys.

\randnotiz{Third—a lover of teaching.} The good Catechist must be a lover of the work of teaching. Realizing that teaching is imparting information, with a purpose, he delights in the labor.

There is a peculiar fascination in the imparting of information. The telling of some new thing is the essence of gossip. The gossiper who tells news, imparts information. He teaches. Whether man or woman, he dearly loves this teaching. He would often rather thus teach than eat. The sensation of telling what the hearer does not know is one of intense delight. And if this imparting of information in ordinary gossip and conversation is a source of pleasure, why should it not be much more so when the very best, the highest, the heavenly truth is imparted, and when it is imparted not for the selfish gratification of the teacher, but for the purpose of making the recipient of the information better and happier for the life that now is and for that which is to come? Yes,

\begin{quote}“To teach the way of life and peace,\\It is a Christlike thing.”\end{quote}

And yet there are pastors who say, or feel if they do not say, “I hate catechising.” May they not have missed their calling? Or have they allowed themselves to become so cold, so worldly, so sluggish,
\quellseite{143}
that they have lost their first love, that their erstwhile zest is gone, that they are like trees that are dying at the top? Gray hairs are here and there upon them, but they know it not. May God help them to be zealous and repent, to strengthen the things that remain, that are ready to die. The good Catechist must love the work of catechising, find joy in it and be sorry when it is over.

\randnotiz{What a good Catechist must know.} Having seen what the good Catechist must be, we inquire next into what he must know. We have said elsewhere that the good minister ought to know everything. But he cannot know it all. Then he should ever strive to learn more and more and to come as near to perfection in knowledge as possible. Since he cannot learn all, he should ever learn what he needs most for the efficient exercise of his high and holy calling. He must be wise in his selection. All this is vitally important for the Catechist. We can merely hint here at the most necessary things for the Catechist to know.

\randnotiz{Bible truth.} We naturally begin with that which must ever form the basis and bulk of all his teaching, which is Bible truth, truth wisely selected, properly ordered and naturally adapted to his pupils. In order thus to select, arrange and adapt, he must know his Bible from beginning to end. He must know the nature, purpose, progress and scope of God’s full revelation.
\quellseite{144}
He must have a clear perspective and understanding of God’s history of redemption.

\randnotiz{Bible facts.} After he has this clear grasp of the Bible as a whole, he must know the content and intent of each book of the Bible. He must know what can be known of the date, the writer and those to or for whom the book was specially written. He must be at home in Biblical Introduction.

Then, too, he must be familiar with Bible geography, antiquities and history. He should be at home in the Biblical theology of the Old and of the New Testament. \randnotiz{Bible doctrines.} He should have a clear understanding and clear convictions of the doctrines of the Bible, i. e., he should know what the teaching is on every subject of doctrine and the relation of these doctrines to each other. In other words, he must have a clear comprehension of a sound, Scriptural system of dogmatics and of ethics as well. He must be a Bible student and scholar, love and constantly study his Bible, live in it and grow in it more and more.

It goes without saying that he must know his Catechism. He must know its history and its contents. He must know it not only verbally and understandingly, but also experimentally. He must know that it is a book for the heart and life. That it teaches not only doctrines, but experiences. That it teaches every doctrine and practice experimentally.
\quellseite{145}
Its explanation of the first Commandment is experimental, as is the explanation of every Commandment. We should fear, love and trust. We repeat, fear is an experience; love is an experience; trust is an experience. This experimental line runs through the whole book. When we come to the fifth and last part, the two questions that close this book of doctrine and devotion, emphasize the words “for you,” as the chief thing in the Sacrament and as bringing the assurance of the forgiveness of sins to him who truly believes these words. This is experience. The good Catechist knows this beautiful and precious characteristic of his Catechism. He lives the experience which the book teaches. He prays his Catechism in his closet and in public. He knows it theoretically, practically and experimentally. \randnotiz{Know the Catechism theoretically, practically, experimentally.} It ought not to be necessary to insist that he know the words of Luther’s parts by heart and can readily repeat them, word for word, line for line and answer for answer. The good Catechist knows how to be simple and clear in his teaching. He uses simple, Saxon, everyday words; words that the dullest can understand and that cannot be misunderstood. He makes his instruction so clear, that all know what he means. He realizes that he is not teaching a class of theological students, or a gathering of ministers or public school teachers, but children. The English Lutheran Catechist must be
\quellseite{146}
\randnotiz{Use simple language.} especially on his guard on this point. Many of his catechumens come from homes where no English, or only the most elementary everyday English is spoken. These pupils are not at home in English literature. They have but a limited vocabulary. Religious terms are especially strange to them. In teaching them as well as in preaching to them all scientific, bookish and class-room terms are to be avoided and only such homely words and phrases are to be used as the simplest use and understand. The good Catechist is a master in the art of simplifying.

\randnotiz{Know how to illustrate.} In connection with clear explanation the good Catechist must know how to use illustrations. The catechumens are not yet beyond the age when they love things more than thoughts. They are still attracted by and to the concrete. It is indeed a question whether any normal person ever outgrows this preference. To make his teaching attractive as well as clear the teacher must know how to illustrate the points he explains. Apt illustrations always help to explain. They convince and satisfy doubt. They also assist the memory. And with all this they interest, attract, please. But they must be well selected. Like our Saviour’s parables, which are still the best illustrations that any teacher ever used, they must be taken from the realm in which the pupil lives and is at home. Illustrations from the domain of
\quellseite{147}
abstract science, from classic art, from musty mythology, from ethnic religion, from the learned lore of fabled antiquity, are all out of place before a catechetical class. In a good illustration the unknown or partially known is made clearly known by showing its likeness to the known. A good illustration throws a flood of light over that which is dark. It thus puts the crown upon the teaching and upon the teacher. Of necessity the illustration must be brief. It dare not overshadow the truth to be illustrated. Sometimes it buries the truth taught so effectively that it is entirely forgotten and the illustration alone is remembered. Such illustrations do more harm than good. They hide the truth, instead of making it clear and prominent. Avoid them. The Catechist who is good at illustration has a great advantage towards making his lessons profitable and permanent and at the same time towards making himself beloved by his class.

\randnotiz{Know how to apply truth.} The good Catechist, again, must know how to apply the truth of the lesson to the heart and life of each learner. He must show not only what the lesson means in itself, but also what it means to each one in the class. Each pupil must be made to see how this particular point of truth is to be used, for warning, for restraint, for encouragement, for comfort or help as the case may be. The lessons of the class should form the conscience of the pupil.
\quellseite{148}
They should be so applied that he will place them before him when in doubt, in temptation, in discouragement or in danger. So will they become the man of his counsel and the guide of his life. Without application the teacher is teaching without a purpose. In other words, he is not teaching at all. True teaching, especially in spiritual things, is the imparting of information with a purpose. Application is the emphasizing of the purpose. It requires a graceful tact and delicate art to be good at application. It dare not run into preaching or exhorting. It may never be angry or harsh. It must always be brief, pointed and plain. Its tone must be kind and loving. Its appeal must be to what is good in the child. Its form should be encouraging and stimulating. It should generally be positive in form. Rarely should it be negative. Character is not built up on “don’ts.” The good Catechist knows how to make good applications.

\randnotiz{Know human nature.} As we have seen, the good Catechist must know human nature. He needs to be a psychologist; to know not only general psychology, but even more does he need to know child-psychology. He must know the child-psychology of the schools. This is very important and helpful for the understanding of the mental makeup, the mental activity and the mental development of child-nature. This study has become a fad in our day. While there is much that is
\quellseite{149}
pure speculation and fancy in the teaching and writing of these specialists, there is also much that is suggestive and helpful to every teacher. The Catechist needs to study and get all the good he can out of these investigations. He will find much that will help him in his catechising.

\randnotiz{Modern child-study defective.} But the Lutheran Catechist will always note a grave defect in modern secular child-study. It studies an unreal child. It postulates an imaginary child. It creates for itself a fanciful child, one that does not exist, a child that has no sinful nature. It is not the offspring of a sin-diseased stock. It does not come from fallen parents. The old truth that

\begin{quote}“In Adam’s fall\\We sinned all.”\end{quote}

is a myth, a fable, which this enlightened twentieth century has long since outgrown.

\randnotiz{Know that sin is in the child.} But the Lutheran Catechist, while he might wish that the modern, rationalistic view were true, knows that is not true. He knows from the Bible, from history, from observation and from experience that the child is shapen in iniquity and conceived in sin, flesh born of the flesh, by nature a child of wrath. He knows the plain teaching of his Bible, from its beginning to its end, concerning sin. He knows that every child is infected by sin. This sin is a realistic part of the child’s nature. It must be
\quellseite{150}
reckoned with. Ignoring it will not banish it. A child-study that leaves it out is untrue. The Lutheran Catechist is a truer psychologist than Froebel, G. Stanley Hall, Professor James, or any of that class. And yet he can learn and should learn much from these men. He must know when and where and how to put in what they leave out. When he does this he is learning a true psychology.

Naturally the Lutheran Catechist, when he considers sin in the child, also considers the nature and efficacy of baptism, as it is so clearly set forth in his Catechism. And so the Lutheran Catechist studies and applies the science of child-psychology and it helps him greatly to be a good Catechist.

\randnotiz{Know the environment of the pupil.} But he needs also to know the environment of each one of his pupils. Environment plays an important part in each life—often a most tragic part. It has an influence and a power that must also be reckoned with. The good Catechist will want to know and will take pains to know from what kind of a home, out of what kind of a life and up from what kind of an atmosphere each one comes. If he remains ignorant of this he will err in his estimates, make serious mistakes in his treatment, fail in making necessary allowance and often do grievous wrong to his pupils. Some of his pupils come from good, pure, refined Christian homes. Others come from sordid, impure, profane, dishonest, coarse and
\quellseite{151}
cruel homes. Others again come from all the varying grades between these two. It is of vital importance that the Catechist know all that he can about this. He must know. He cannot be fair to his class without it. He cannot be just without it. He cannot properly either explain or apply the truth without it. But with this knowledge he will have a magic key to open the different hearts and minds. He will have a search-light to illumine the dark places in the hearts and habits. He will be able to be more earnest, more tender, more effective. He will be able to throw himself into each life more heartily and more enthusiastically. The good Catechist must know the environment of each and of all of his class.

\randnotiz{Know pedagogics.} And finally, the good Catechist must know the history, the science and art of pedagogics. Before his class he is a pedagogue. How can he be a good one, if he knows nothing about teaching? If the State will not permit anyone to teach even the most elementary backwoods school without first giving evidence that he has studied and knows pedagogy, why should the Church send out men to catechise her children without this knowledge? Have not our theological seminaries made a grievous mistake in not having a good course of pedagogics as a necessary part of their curricula? Yes, the good Catechist must know pedagogy. He must know how others
\quellseite{152}
have taught; what has been the secret of the successful ones; what the causes of the failures; what belongs to real teaching; to the personality of the teacher; to his demeanor before the class; to his management of the class; to his treatment of the various ages, characters, abilities and attitudes of his pupils; what belongs to a cheerful, yet firm and effective discipline; to explaining, impressing, clinching, applying the lesson. All these and many other needed things can be learned from a good course in pedagogy. The good Catechist needs to know pedagogics.

\quellseite{153}
\chapter[Chapter XIII. Gathering His Class]{Chapter XIII. The Catechist Gathering His Class}

CHAPTER XIII.

THE CATECHIST GATHERING HIS CLASS.

The properly trained, equipped and qualified Catechist is ready for his work. He is eager for it. He looks forward to every new class with glad anticipation. How shall he get his class? Whom ought he to gather in for instruction?

\randnotiz{Gathers the baptized children.} As a matter of course he wants and expects the baptized children of the Church. They have become members of Christ’s Church by His own initial Sacrament. They are in a peculiar sense His lambs. They are to be fed, nurtured, developed and trained so they not only may be and remain His, live under Him in His Kingdom, but also serve Him in everlasting righteousness, innocence and blessedness. For this they have been redeemed. For this redemption has been applied in holy baptism. For this they are called, gathered, enlightened and sanctified: That they might be His and also that they might serve. They are to be the future Church in which men are to find forgiveness and salvation and by which these blessings are to be brought to others also.

The good Shepherd is not satisfied with the flock that has been gathered and saved. He has a will
\quellseite{154}
and desire for the ungathered and unsaved. “Other sheep I have which are not of this fold. Them also I must bring, and they shall hear my voice.” By the hearing of His voice they are to be gathered in. That voice is sounded out by His under-shepherds, His ministers, His Catechists. If there are other sheep, they have lambs. These are outside. Them also must He bring, through the under-shepherd.

\randnotiz{Gathers the neglected.} The good Catechist, therefore, is not satisfied when he has gathered the baptized children of his congregation. He has an earnest care, a loving and longing solicitation for the other sheep and lambs. Them also he would bring. He is constantly seeking the lost and erring. He has a scent for every unshepherded lamb in his reach. In his canvass for his own children, whom he would gather into his class, he is on the lookout, inquiring and searching for neglected children and youth and older ones.

How earnestly he seeks and pleads that he may get them into his class, have the privilege of instructing them and making them wise unto salvation. He who has no care for, and makes no effort for, the neglected is not a true under-shepherd, is not doing the work of an evangelist, is not making full proof of his ministry, is not a good Catechist.

As the good Catechist goes from house to house seeking out and persuading all such to come, he makes it clear that he promises no one Confirmation.
\quellseite{155}
\randnotiz{Promises no one Confirmation.} He is going to confirm only those who desire it and of whose fitness he is satisfied. In what that fitness consists we shall consider in a future chapter. But how could he know beforehand who is going to be fit? The instruction and the manner of receiving it are to bring about the fitness. When the course is finished and the final tests have been made, then and not till then does he promise Confirmation. Let him make this clear to all parents and to all outsiders, both parents and those who are invited to instruction, that he promises Confirmation to no one at this end of the line. No parent and no catechumen has a right to demand Confirmation. The Catechist who makes this so clear that it cannot be misunderstood, will save himself from some serious embarrassment in the future. By neglecting this important practical point many a Catechist has brought upon himself serious trouble.

\randnotiz{Attendance does not commit to Confirmation.} On the other hand, in the canvass for members of his class, the good Catechist also makes it clear that no one who comes to instruction thereby commits himself to be confirmed. Should some not want to be confirmed at the end of the course the Catechist does not regret that he instructed them. He is thankful that he had the privilege. He has sown the seed. It is not necessarily dead because it does not bear fruit at once. It may spring up and bear fruit in future days. He dismisses such catechumens
\quellseite{156}
at the end of the course with his prayers, his blessing, his most loving counsels. He has cast his bread upon the waters. He shall find it, or another may find it, after many days. Perhaps in another church, in another clime, at another time, the fruit will ripen.

In his canvass the good Catechist will gather all the information he can as to the families from which the children come. We have seen the important influence of environment and the need for the Catechist to fully understand it.

\randnotiz{Seeks to enlist aid of parents.} He will also do his best to enlist the interest of the parents in an earnest heart-to-heart talk with father and mother on the seriousness of proper catechising, on the pastor’s need of the parents’ interest, prayer and co-operation. They must be made to understand that if they take no interest it is likely that the children will take none. It should be earnestly impressed upon the parents that it rests with them whether the children will come regularly and be in their places promptly. It is their responsibility, their bounden duty, to take an interest in the lessons, to encourage regular, daily preparation, to inquire into that preparation and to assist in it. Impress it upon them that no Catechist can do his best without such interest and co-operation. The Catechist can counsel, admonish and exhort that at least half an hour be given every day to earnest preparation of
\quellseite{157}
the next lesson. But he cannot enforce this very necessary practice. The parents can. They can and should use inducement, assistance and parental authority. They must understand that there can be no progress and improvement such as the Catechist desires, where there is no home preparation, or at best, a hasty skimming of the lesson immediately before leaving for, or on the way to class. The parents as well as the children must understand that the right result cannot be reached unless there is a part of each day scrupulously given up for such preparation. The parents must also be reasoned and pleaded with to give the pastor their earnest and prayerful assistance in bringing about and increasing heart-fitness.

\randnotiz{Limitations of parental authority.} The parents must also understand where their authority ends. They have not only the right, but the duty to insist that the children prepare and attend regularly. They must assist the pastor in compelling proper behavior. But all this in the true Christian spirit. They dare not provoke their children unto wrath. If they are the right kind of Christian parents they will use only the compulsion of love. Religion cannot be commanded or beaten in with the rod. No authority can force religion into the child. Even God has so ordained that He cannot do it. He pleads, calls, influences, exhorts and commands. But He does not force. The parents dare
\quellseite{158}
not command their children or demand of the pastor that they must be confirmed. The baptismal vow was not consciously voluntary on the part of the child. The Confirmation vow, as we have already seen, must be voluntary. Its very meaning is gone, if it is not taken intelligently and with free will. It becomes a mockery, a pious fraud. The good Catechist will see to it that the parents understand that here their authority ends.

\randnotiz{Special sermon in advance.} It is well for the pastor also to preach a special sermon every year several weeks before he starts each class. In this sermon he should set forth the need of religious instruction, especially in our day, the Bible idea, the Church idea, and especially the idea and view of our own Church. He should show the importance and value of such a course for the child, for the family, for the Church and for the State. He should impress the fact that no other church has such a high appreciation of the child, such a clear understanding of child-nature and its needs, such an attractive and effective scheme in its view of infant baptism and Christian nurture for supplying the needs as has our Church. He should plead for the interest, the co-operation and the prayers of the congregation, especially of office-bearers and parents, in making the most of the good old custom of catechisation. Thus the sermon and the canvass should support and supplement each
\quellseite{159}
other. The sermon should precede the canvass and the canvass should clinch and carry out the intent of the sermon.

Before he announces the starting of a new course of instruction the pastor must determine in his own mind the time and place for the meetings.

The place, as a rule,\randnotiz{Place.} is the most suitable, available room in the church. In exceptional cases it may be better to meet elsewhere. Small classes may meet in the pastor’s study or in some other private house.

The time is not always\randnotiz{Time.} so easily determined. The time of year is usually from autumn till Easter or Pentecost. Better the latter, as it gives more time, and all the time is needed.\footnote{On the sparsely settled prairies of the Northwest, with its killing cold and burying blizzards, it is, of course, unadvisable to have catechisation at this time of the year. Common sense teaches that the time for instruction is summer and the time for Confirmation late autumn.}

As to the day or days of the week, this will depend on the opportunity and convenience of the class as a whole. But it dare not be left to the choice of the class. This would cause division and dissatisfaction. The Catechist must carefully make up his mind and then announce the time. The attendance of public school must be considered. We do not believe that the State schools will ever give up part of their time for the Church.\footnote{See a full discussion of such a proposal in “Religious Education and the Public School,” by the Rev. G. U. Wenner, D. D.} Neither can
\quellseite{160}
it be expected that our catechumens lose part of every week’s school time and so lose their standing in their classes. We must reckon with the public school as it is. Our classes must meet when our children can be present, without being absent from public school. Saturday forenoon or afternoon will generally be best.

Frequently the Catechist will find that all who would and should come cannot come at the same time. The children cannot come at night. Others are employed during the day and can come only at night. For these a night class will have to be arranged. The zealous Catechist will be glad to do the extra work which this entails.

Occasionally he will find some who cannot attend class at all. These he will have to instruct privately.\randnotiz{Private instruction.} They may be fathers or mothers, too tired to come or ashamed to take their places in class with boys and girls. An unnecessary cross ought not to be laid on them. The pastor should gladly go to them. On account of age and mental habit it may be next to impossible for them to memorize. There are pastors who have so little common sense that they refuse to receive such persons into the congregation, unless they commit and recite the Catechism. Like the Pharisees of old they would bind on heavy burdens which neither they nor their children would be willing to bear.

\quellseite{161}
The good Catechist knows how to make allowance. He will be thankful if he can get such people to have an intelligent comprehension of the most essential teachings of the Church concerning sin, grace and the means of grace. They can often get more good from reading a good explanation than from following the questions and answers of the Catechism. Let the pastor put the right book into their hands.\footnote{It was partly to reach and assist such people that the writer of this wrote “The Way of Salvation in the Lutheran Church.”} And so the good Catechist suits his own time to the best convenience of all, who are willing to be instructed, and, in the good sense, gladly becomes all things to all men, if by any means he may save some. Such a Catechist will always find people to instruct and will find joy in his instruction. And surely his labor will not be in vain in the Lord. He will have souls for his hire.

When he announces his class, he will also announce his text-book. What explanation of Luther’s Catechism shall he use?\randnotiz{What explanation?} Their name is legion and still there are more to follow.

The good Catechist will want on his book shelf, for his class preparation and private use, Spener’s Explanation, Loehe’s, Dietrich’s, Stump’s, Reu’s. But we do not advise any of these for class use, though we would make no law on the subject. The Augustana Synod’s English Explanation is excellent.
\quellseite{162}
for class use. Among Norwegians there is nothing to be preferred to Pontoppidan’s Explanation, which is based on Spener’s. The pity is that we do not have a translation into simple, Saxon English. Bishop Laache’s abbreviation of Pontoppidan is also good. The English translation is much better than that of Pontoppidan. For ourselves we have found our best satisfaction in late years in the use of Dr. Trabert’s clear, terse, idiomatic questions and answers. But let everyone be fully persuaded in his own mind and use that explanation which will bring about the best results in his work.

Let the Catechist always have a good supply of his chosen text-book on hand. Let him announce that the books can be secured from him. Let him sell them at cost to his catechumens. It is a shame to try to make gain for himself in supplying Catechisms or any other church books to his people. Should there be some so poor that it would be a burden for them to buy, let the pastor donate books to them. He is everlastingly preaching to his people to give. He ought to be the most liberal giver of all.

And so he has gathered his class. All things are ready for them. He eagerly looks forward to the first meeting. The conduct of the first meeting we shall consider in the next chapter.

\quellseite{163}
\chapter[Chapter XIV. The First Meeting]{Chapter XIV. The Catechist Before His Class. The First Meeting}


The hour for the first meeting\randnotiz{Things that count.} of the new class comes. The meeting place must be ready and open in good time. The place should be clean, comfortable, cheerful. The Catechist should be on hand not only at this meeting, but at every meeting fifteen or thirty minutes before the time set for beginning. The time of waiting need not be wasted. It will be valuable for final preparation. He should greet everyone cheerfully when he arrives, see that all are put at their ease, that wet or cold feet are attended to, and that no disorder or mischief be started before the opening. It is easier to prevent than to quell disorder.

The Catechist should be at his best. First impressions count. A good impression made at the first meeting will count through the whole course. He must know how to be cheerful without descending to the undignified, the irreverent, the trifling. He should never allow himself to have a forbidding mien or tone. He should never be cross or harsh. He should never scold. This applies to every meeting, but is specially important for the first.

\quellseite{164}
Let him begin promptly on time. Impress the fact that you will always begin on time and then live up to the rule. Call the class to order.\randnotiz{How to begin.} Wait till there is perfect order. If one is whispering, say quietly, “John, or Mary, we are waiting.” In a quiet word remind them that all are to pray with you. Then say, “Let all arise and pray.” Pray a short, free, earnest, direct prayer, giving thanks for the privilege of studying together what alone can make wise unto salvation; praying for assistance and guidance for this particular lesson, for the class as a whole and for each member; for any that may be sick, in temptation, danger or distress; for the forgiveness of sin and for deliverance from sin for this life and for the next. Should there be whispering during prayer, stop, look at the whisperer, say nothing, and go on with your prayer. After the prayer say quietly how sorry you are that someone forgot and was irreverent. Don’t scold.

If you can lead the singing and have a few good voices in the class, it is well to sing, not more than two or three stanzas from a familiar hymn. Better not sing than sing a solo. It would be well to insist that each member of the class own a Church Book as well as a Bible and bring both to class along with the Catechism each time. If, after the prayer, a few minutes, not more than five or eight, can be taken regularly for the explanation of the Church Book,
\quellseite{165}
its table of contents, the Common Service and some of the best hymns, it will be a blessing to the class. Then make up your roll. Put down the names\randnotiz{Keep a roll.} in alphabetical order. Have a blank book which is convenient for marking absences. You need not mark those present. To be present is taken for granted, to be absent needs to be noted. Have a kind word, a word of welcome, for everyone who has not yet had it. If there is sickness in a family, inquire kindly when you come to the name of a member of that home. Let all feel that the pastor is interested in them, is their friend and the friend of their families.


After the roll is made up, inquire how many have Bibles of their own. See that each one gets a good one. If any are too poor to secure Bibles and Church Books of their own, and if the pastor cannot donate all that are needed, he should see that the congregation promptly supply the need. The Catechist should also take a part of each session, as much as is commended for the Church Book, to teach the class the order of the books of the Bible. This will require pedagogic skill, constant drill and repetition. Teach the books of the Pentateuch\randnotiz{Three needed books.} first, then the historical books, then the poetical, the prophetical, the Gospels and Acts, the Epistles. This will require constant review. It is important. Catechumens must be able to handle their Bibles. Test them often by having
\quellseite{166}
them look up passages. This is for every meeting. To come back to the first. After the Catechist has ascertained who has and who does not have the three needed books and arranged to supply the lack, he will give his class a well-prepared, free and informal explanatory talk.


He will explain the purpose\randnotiz{Explanatory talk.} of the course of study before them. He will show them that they are to learn what God wants them to know, that they are not only to know what God teaches, but they are to learn to live up to it. This course is expected to be the beginning of an earnest endeavor to grow into a stronger and closer discipleship with Christ. They are to grow in grace as they grow in this knowledge. They are to have an ever-increasing desire and determination to love Christ, to live for Christ and to labor for Christ. He reminds them that to make the course profitable for all he must expect regularity, promptness and valid excuses for absence or tardiness.

He must tell them very firmly and impressively, though kindly, that in class he will insist on attention and order, that he cannot and will not teach without order, that if there are any who are not willing to behave and obey as he desires they cannot be in the class. These necessary principles he must carry out without hesitancy or variation.

He must make clear to the members of his class
\quellseite{167}
what he expects from them at home. He must urge upon them\randnotiz{Home devotions and preparation.} the importance and necessity of daily devotional Bible reading, of forming regular habits of prayer, of praying for God’s blessing on the study of the Catechism and the attendance of class. He expects from each one diligent, daily preparation of the lessons assigned. Let him impress the fact that they cannot play or bluff their way through, that it means study, preparation, repetition, self-quiz, review; that they must rigidly set apart a portion of each day for these lessons and be in earnest in improving them, that no one can take the course by depending on a hurried skimming over of the lesson just before starting for class. He cannot make this point too impressive.

Again, let him enforce the point\randnotiz{Purpose of course.}, which will come up in lesson after lesson, that their purpose should not be to get through and be confirmed, that Confirmation without earnest and fruitful preparation would be an empty ceremony, and that he promises Confirmation to no one before the course is ended and the final test is had. They are to take this course that they may learn God’s truth, that they may understand it, that they may have an ever-deepening experience of its power and preciousness.

Assure the class that you anticipate pleasant meetings, that they will be pleasant for everyone in the class who will do his part faithfully. Dwell on the
\quellseite{168}
pleasures of the course. Speak of the value of the course\randnotiz{Pleasure and value of course.} for the whole future, show how it has always brought blessing where rightly used. That no one had ever regretted it, but many lament not having had the privilege. Give examples. The writer of this used to tell his classes of a blind woman, whom he knew, who found light and comfort in her darkness by repeating to herself Luther’s Catechism. Of another woman who had become so childish in her old age that she did not know her own children, and how the writer once overheard her going over the shorter Catechism—she was a devout Presbyterian—beginning with the first question and answer, then the second, and so on. Poor soul! she didn’t know her own children and could get no comfort from them, but she knew the Catechism that she had learned three-quarters of a century ago and found comfort and peace in its old lessons. It is said of Andreae, one of the framers of the Formula of Concord, that when he was very old and childish he imagined himself a boy again. He would call to the companions of his boyhood to come and have a game of ball. Again he would say, “Boys, do we know our Catechism lesson?” Then he would repeat question and answer from any part of the book, most frequently from Luther’s Explanation of the Apostles’ Creed. And so the treasures of the Catechism, when engraved on the memory, remain as a
\quellseite{169}
sweet solace and stay when earthly comforts fail. Examples encourage the class.

And, finally, impress upon the class\randnotiz{Faithful attendance.} that you expect everyone to be faithful and regular in attending the Church service, the Sunday school and the Luther League. Show them that if they neglect these services they show a lack of taste and interest in holy things, which does not recommend them for Confirmation; that if they do not habitually attend now, they will not be likely to do so after Confirmation, and that you should not want to confirm anyone who would not expect to attend all these services after Confirmation.

Have the catechumens sit together in a front seat in church. If the class is large, appoint a member to keep an attendance roll at church. Inquire into this roll at your class gatherings and speak of the importance of church attendance frequently. Assign the lesson for the next time.\randnotiz{Directions for next lesson.} Explain what you want them to memorize and show how you want them to study the rest. It might be well to assign weekly portions for the daily Bible readings. Do not make them too long. Do not assign too much anywhere. Not how much, but how well, is a good pedagogic principle.

After another brief prayer, in which you pray for a blessing on the hour spent together and on the lessons received, dismiss the class with a kindly goodbye.

\quellseite{170}
\chapter[Chapter XV. An Ordinary Regular Meeting]{Chapter XV. The Catechist Before His Class. An Ordinary Regular Meeting}

On the opening and closing devotions the principles laid down in the former chapter will apply to all ordinary meetings of the class. The principles are the same, the method will vary with the varying conditions and needs. The good Catechist, like the good preacher, will always be timely. He will know how to apply unchanging principles to the conditions and needs of the hour.

We would again remind the Catechist\randnotiz{Conduct of regular meetings.} that he is to have himself well in hand, is to be prepared in mind and heart, is to be there in time to greet each one as he comes and to prevent incipient, disturbing forces, is to be calm, kind and cheerful. It is well to call the roll before and not after the lesson. This custom is an incentive to promptness. In the roll-call let him ask after absentees, sickness in families, and other matters that call for interest and sympathy.

When perfect order is secured after the devotions and the brief opening drill on contents and use of Bible and Church Book, as recommended above, let him have a very brief and stimulating review\randnotiz{Review.} of the main truth and trend of the last lesson. Here a few,
\quellseite{171}
a very few apt questions, intended to refresh the memory and clinch the truth are in place. By a skillful review of this kind he will link the last lesson to the one before the class. He will thus put into practice that pedagogic principle called apperception, so much emphasized of late by the followers of Herbart.

Now let him go forward with the lesson in hand. Let there be no careless delays, no idle pauses which are sure to encourage listlessness and disorder. He must get attention to start with. He must hold the attention. If he lets it lag he loses his grip on his class. He must keep their minds interested and intent on the lesson.

Much useless discussion\randnotiz{Attention and interest.} has been had on the question whether attention must precede interest, or interest attention, and on the question of the relation of the one to the other. Labored argument has been had in favor of the necessary precedence of the one or the other. Taking a common sense view of the matter, we believe that such argument is useless.

Attention and interest go together. They cannot be separated. You cannot have the one for any length of time without the other. A demanded and forced attention without interest is barren of results. When you call for order and wait on order you get attention. Now at once give them something to attend to. Get them interested. Attention opens the door
\quellseite{172}
for interest and interest holds fast the attention. Profitable attention at once merges into interest. Attention may be, and in children doubtless is at first an effort. The conscious effort of the will is in it. When it merges into interest it becomes unconscious and in that sense involuntary. Then it is fruitful. It is this interest that the Catechist wants to have and to hold. To this end he must keep all busy. How can he do it?

Having first secured order, held it through roll-call, devotions, Bible and Church Book drill, he must go rapidly into the lesson for the day. He will ask first a few questions\randnotiz{Questions—method of.} that are not in the book, as to what the last lesson was about, what was of special interest in it and such like. He will repeat his question clearly and have it understood that the one named is to answer. He is never to question by rote or in the order of the sitting. Such a custom is an encouragement to inattention and cheating. But after he has asked his question he will call on one here, one there, one elsewhere, as he names him, to answer. His eyes will constantly take in the whole. He will notice every disorder, distraction, mischief, whispering, noise. He will ask the inattentive and disorderly one for an answer. That one may say, “What is the question?” The teacher says, “I asked the question clearly; if you had paid attention you would have heard.” This practice of asking the
\quellseite{173}
question and then quickly calling on the inattentive one for an answer is wonderfully effective towards maintaining general interest and good order. For other reasons also the practice of skipping about for answers is commendable.

The Catechist needs to be guarded also against favoring the brilliant and forward ones in the class. They often need to be held back and repressed, lest they monopolize the time of the teacher and discourage the rest of the class. Let them understand that they are not to answer until they are named. Let the Catechist rather favor the timid and the dull.

Ask them easy questions, encourage and help them, show them that they can answer. They will be glad to do so. Ask the harder questions of the brighter ones. To whomsoever much is given, of him also much will be required. Remember also that some have encouragement and help at home, others have none. Make allowance.

After a few connecting and introductory questions of your own, turn to the questions of the book. To these link your own questions. Your own should further elucidate, develop, illustrate and apply those of the book.

The need, the nature and the form\randnotiz{Sequence of questions.} of the questions to be asked we have considered above. It is not necessary that we repeat here what we said in Chapter VII. We would again impress the importance
\quellseite{174}
of the sequence of the questions. They should follow each other in natural and logical order. The negative should come before the positive. False conceptions and errors should be disposed of first. When the false has been put out of the way, the true should be implanted, made clear, impressed and applied.

Here we might also inquire\randnotiz{Socratic questions.} into the place of the so-called Socratic questions in the catechetical class. We certainly cannot use them in the sense and for the purpose that they served Socrates. He used them with telling effect for his own purpose. He dealt with natural, innate truth and its development and application. The truth whose germs lie hidden in the mind of the natural man, which the natural mind can recognize, grasp, expand and apply, this was the material on which Socrates brought his method to bear. He wanted his pupils to learn to use the powers of their own minds correctly. He wanted them to think cautiously and clearly, to reason correctly, to conclude rightly. Dealing with such truth as is indigenous to the human soul and with such powers as belong to every normal mind, Socrates rendered a most valuable service to his own and to every succeeding age, by his original, unique and far-reaching method. Every Catechist, every preacher, should have a clear comprehension of the Socratic method of teaching.

\quellseite{175}
But the Catechist is not dealing with the natural alone. He is dealing with the supernatural. The truth which he teaches\randnotiz{Divine, saving truth.} is not native to the natural man. “The natural man receiveth not the things of the Spirit of God, for they are foolishness unto him, neither can he know them, for they are spiritually discerned.” Read 2 Cor. 2. But the things which the natural man could neither originate out of himself or of nature, God has revealed. It is this revealed truth—truth which neither the philosophers, nor sages, nor princes of this world knew—with which the Catechist deals. Since this divine, spiritual, saving truth is not in man by nature, it cannot be drawn out of or developed from his inner consciousness. It must be put into man before he can have it. And this is the Catechist’s business. He is to take the things which God makes known and offers in His Word and give them to his pupils. He is to so implant, instill, lodge and impress God’s truth that the pupil receives, accepts and assimilates it into his very being. He puts it into his mind, into his heart, into his life.

After he has first imparted the truth in the memory of the pupil, then there is something in the mind which was not there before. On this something he now goes to work. On this he can bring his Socratic questions to bear. If he knows how to rightly use them, he can, through them, make the pupil more
\quellseite{176}
fully understand, appropriate and apply them. Through skillful Socratic questioning the pupil gets a clearer and firmer grasp of the truth and of its bearing. He gets visions of its attractiveness, its usefulness, its power. It becomes more and more interesting. Its many-sided relation to other truth is seen. It becomes more and more precious. It is loved for what it is and for what it does.

In using the Socratic questions\randnotiz{How to use Socratic method.} let the same method be followed that is recommended above for the other questions. Ask the question first. Then name the person from whom you want an answer. Ask the easier ones from the less able and the harder ones from the more ready and able. The Socratic method can be used effectively to curb the too forward and conceited ones. Use it as Socrates did, to show them their ignorance and teach them the needed humility. If the pupil asked cannot answer correctly, ask another, and if after several attempts you fail to get what you want, ask how many in the class can answer. Then select one from the raised hands and let it be understood that no one is to answer till named. After you fail to get a satisfactory answer from an individual, you may say, “Will anyone answer?” Repress firmly the would-be smart answers, all answers that are intended to raise a laugh. When you cannot elicit the right answer, then give yours. Give it so clearly, that it
\quellseite{177}
% Continues the preceding page.
cannot be misunderstood. Then quiz again to get the right answer from those who failed or gave the wrong one. Do not pass the point until the class has understood and appropriated it. When thus properly used, the Socratic method can be effectively used to make clear what has been memorized and to so lodge and clinch the truth that it will never be forgotten. This we believe to be the place and use of Socratic questions in catechising.

At the close of the lesson let the Catechist give a brief résumé of what has been studied in this hour and how it should help all to be better Christians. Then give a résumé of the next lesson, pointing out the things of special interest and importance.

\randnotiz{Dismissal.} Close as indicated with a short, free prayer, praying for a blessing on the time spent together, on the lesson learned and commending each and all to the guardian care of the heavenly Father, to the redeeming love of Christ and to the sanctifying grace of the Holy Spirit. End with the Lord’s Prayer in unison. Dismiss with a pleasant word and the hope of all coming next time.

\quellseite{178}
\chapter[Chapter XVI. General Suggestions]{Chapter XVI. The Catechist Before His Class. General Suggestions}

Many a Catechist’s work and many a catechetical hour fails of its proper fruitage because of seemingly small mistakes in the conduct of the work. A chapter of practical suggestions, therefore, will not be amiss.

\randnotiz{Parental attendance.} 1. We have shown that it is of vital importance to have the interest and co-operation of the parents. They ought to be deeply interested in what their children are being taught, how they are being taught and how they receive the teaching. Let the pastor, therefore, not only invite, but earnestly urge the parents, both publicly and privately, to attend. They will thus know what the training of their children is, but they will also carry home important and helpful lessons for their own spiritual life. Not all can come. Some can come now and then, others do not care to come. But the coming of some, even occasionally, will do good. It will do good to them. It will do good to the class and it will do much good to the teacher. He will be encouraged. He will prepare himself more thoroughly, he will do better teaching.

\quellseite{179}
% Continues the preceding page.
\randnotiz{Length of session.} 2. As to the length of each session. The good Catechist wants all the time he can get. He feels that his time is all too short. And there is too often a senseless fear of making the hour too long. The same children will sit in the public school for five or six hours daily for five days in the week. They think they are terribly abused if they are kept for two and one half or three hours in catechetical class once or twice a week. This feeling ought to be discouraged in the start. Parents and pupils should be given to understand that this course of instruction is a serious business, and that it means work, earnest, hard and continuous work. A period of two hours for the daylight classes is short enough. Better get as nearly three hours as possible. The evening classes will have to be shortened according to circumstances. Let the Catechist not be afraid to ask for service and for sacrifice. People prize that which costs something in labor or in means. They esteem lightly that which they get too easily.

\randnotiz{Good order.} 3. The Catechist cannot do good work unless he has good order in his class. He must insist on perfect order. This is specially important during his first meetings with his class. If he once loses his grip and the class becomes disorderly, it will be hard for him to regain it. If he secures perfect order during his first hour he can maintain it.

Let him always be in his place ahead of the first
\quellseite{180}
% Continues the preceding page.
arrivals, speak a kindly word to each one and have them all take their places and get busy with their books. If they gather ten minutes before he appears they will start all sorts of mischief and it will be doubly hard for him to compose and hold them. He should never scold. If he shows anger he loses his grip. The class will not respect him. He has lowered himself, and henceforth is weak. The majesty of calmness is needed. It compels respect. Firmness—serene and kindly firmness—wins. A silent, searching look at the restless, whispering or laughing one will generally suffice. Sometimes a calm, severe word of reproof may be necessary. But don’t scold.

Should there be disorder during prayer, stop and look at the offender, and then finish the prayer. A word as to the sin of irreverence, after or before the devotions, may be in place. During the course of the lesson the Catechist should have his eye on all the class, notice every inattention and disorder, and, as already suggested, spring his questions promiscuously, especially on the inattentive. Let him never forget that he cannot do good teaching without good order.

\randnotiz{The Bible.} 4. We have already spoken of the use of the Bible. Every pupil should have one and bring it to class every time. Home readings should be regularly assigned, and all should be regularly examined
\quellseite{181}
% Continues the preceding page.
on what they have read. The drill on the order of the books of the Bible should run through the whole course. Their knowledge of the order and location of the books should be tested at every lesson by having the pupils look up the reference passages.

\randnotiz{More needed than mere knowing.} 5. Remind the pupils at every lesson that mere knowing is not enough. Make clear and insist that the heart and will must be right. Get them to clearly understand what you mean by the heart being right, show them again what the sinfulness and damnableness of sin is, what hatred of and sorrow for sin mean. Use illustrations of other sorrow and hatred. Portray the experience of these emotions. Make them real and show what real sorrow for sin is, what real hatred of sin is, what real desire for forgiveness and freedom from sin is, what heart-felt confession is, what real prayer for a new heart and a right spirit is. Rest not till everyone understands what real, evangelical heart-religion means. Forget not at a single meeting that the instruction is but a means, that the end is fear, trust in and love to God in Christ.

6. How much shall the catechumen commit to memory? How much shall the Catechist insist on? What is the average? What the minimum?

These questions trouble the Catechist. They cannot be answered categorically. In general we still believe the old-fashioned doctrine, that the more
\quellseite{182}
% Continues the preceding page.
divine truth the child can store away in the memory the better. But all are not equally gifted. \randnotiz{Memorizing.} We do not believe what some teach, viz., that all memories are equally good and the difference in attainment lies in the using. We believe that there are strong and weak memories. What is easy for one child to commit is very difficult for another, and may be impossible for a third. The writer of this has frequently found children for whom it seemed a physical impossibility to commit Luther’s Explanation of the first article of the Apostles’ Creed.

Then, too, not all have the same amount of time at their disposal, nor the same encouragement and help at home. For these things, also, the wise Catechist will make allowance.

Our advice is that the average catechumen commit to memory the five parts that Luther wrote, with as many Scripture passages as possible—always the most important first. And when one is found who for some of the above reasons cannot commit even this minimum, let the Catechist do the best he can to get such an one to understand the meaning of what cannot be memorized. Let such an one not be scolded, but kindly encouraged. Dullness and lack of opportunity will not debar from the kingdom of God, and should not debar from the communion of the Church. This applies also to such as are too old to memorize. To this we have referred above.

\quellseite{183}
% Continues the preceding page.
7. \randnotiz{A closing review.} When it comes to the closing of the course, we advise one or more hours for a review of the textbook. In this review—to which the parents and members of the church council should be invited—the committed text should be repeated and impressed, together with such verbal explanations as have been instilled into the memory. This is the intellectual review. Let the Catechist be careful that he do not favor the bright or forward pupil. Let him be fair and gentle toward the dull and poorly favored one. Let him be well prepared for this review and make it impressive not only to the class, but also to the visitors. To the skillful Catechist this will be an evangelistic opportunity. Over the heads of his pupils he will speak many a word of admonition, of warning and of earnest reminder to those present, who have forgotten many of the old lessons.

8. \randnotiz{A Socratic review.} At least one meeting of the class should be devoted to a Socratic review. To this, church council, parents and friends should be urgently invited. This is probably the most important, as well as the most difficult lesson in the whole course. For this he needs special prayer and painstaking preparation. In this review he is to probe not for memorized truth, but for assimilated and digested truth. The pupil is to tell, in his own words, what he knows,
\quellseite{184}
% Continues the preceding page.
understands and believes on each point on which he is questioned. In this review the teacher will see in how far he has succeeded in having the pupil make the truth taught his own possession, so really his own that he can at any time give it to others in his own words, discuss and defend it. Here the teacher will discover where he has been weak, where he has failed, and how much he still lacks of being a perfect Catechist. This Socratic review will always humble him and drive him to repentance and earnest effort for improvement.

In this review the teacher will find out how much his pupils can tell of the various uses of the law: what it means for the unregenerate and what for the regenerate, and how it should be used in the daily life of the child of God; the pupil will tell what sin is and what its fruits are, where it is unpardoned and unchecked by grace and godly fear and watchfulness; he will be able to tell what he knows of the nature of God, of His attributes, His work in creation and in providence, of His revelation of Himself, and of His will in His Word; he will tell of Christ, of His person, His work, His advent and its meaning to the unsaved and to the saved; he will explain what he knows of the Spirit, of His person, His operations, His means and methods, the subjective influences of these operations, the results of
\quellseite{185}
% Continues the preceding page.
accepting and of rejecting these divine influences. And so he can go on and tell of prayer, its nature, its conditions, its relations to things temporal and to things spiritual, its application to each one of the seven petitions. The Catechist will see whether they are clear on prayer, whether they have Reformed or fanatical ideas, which are productive of so much skepticism. The teacher will thus go on and see how clear their own ideas are on confession, on baptism, on the difference between Lutheran, Baptist and general Reformed doctrine, on the Lord’s Supper, what they know of the false views, of the Romish and the Reformed over against the true view.

\randnotiz{An opportunity to evangelize others.} And here, again, is the pastor’s opportunity to instruct, to evangelize the visitors. He may allow himself a wide range of application. He may dwell on whatever is most important to his own people. He may emphasize pet sins, Church going, Church support, danger of particular phases of worldliness, or whatever is most needed. The Socratic review is a difficult one. But it can be made a most blessed one for the Catechist, the class and the visitors, who should be many. There ought to be ample time for this review. The true Catechist will insist, as we have insisted all through this book, on knowing first, “What does this mean in itself?” and second, “What does it mean to me?” and third, “How can I
\quellseite{186}
% Continues the preceding page.
experience the full benefit of what it means to me?” Such a review is one of the crowns of the catechetical course. Blessed is the pastor who knows how to get the full benefit out of it all.

\quellseite{187}
\chapter[Chapter XVII. Whom He Will Confirm]{Chapter XVII. The Catechist Deciding Whom He Will Confirm}

\randnotiz{Whom to confirm.} We have discussed the solemn rite of Confirmation. We have seen that it is not a divine institution, and, therefore, not necessary to salvation. It is of vastly more importance to instruct and fit for Confirmation than to confirm. Knowing this, the good Catechist is much more eager to teach the truth which sanctifies and to impart the knowledge which is life eternal than to lay on hands in Confirmation. He has as yet promised Confirmation to no one. He is, therefore, perfectly free to confirm or not to confirm. But he must decide which he will do. Though he has been careful not to promise, yet many parents and pupils desire Confirmation. The Catechist must decide whom he shall confirm. How shall he decide? On what grounds shall he base his determination? What are the true tests for Confirmation? How can the tests be made? What shall the Catechist do so as to make no mistake? What does he want in his catechumens?

He wants, first of all, intellectual fitness. This has been insisted on all through the book. There is to be a clear apprehension and comprehension of
\quellseite{188}
% Continues the preceding page.
the great foundation truths of God’s Word. The “form of sound words” as contained in the Catechism is to be an abiding content of the memory. The more firmly the sound words are lodged in the catechumen, the better does the Catechist like it. But, as already shown, he cannot and does not expect the same memory work and fruit in all. Still, he wants all he can get.

\randnotiz{Clear understanding.} He wants a clear understanding of the truth studied and memorized. He wants the pupil to have a clear insight of what he has learned. He wants the pupil to be able to tell and explain in his own words the facts and truths he has learned. The pupil is to be able and ready to state and to defend the true teaching. He is to be able to set forth the difference between the true and the various phases of false teaching on any important point of doctrine. He is to be ready always to tell what he has been taught to believe and to give reasons for his belief. Thus the Catechist wants intellectual possession of truth, clear discrimination between truth and error, ability to clearly state and distinguish and also ability to defend his acquirements and to confute the gainsayer.

\randnotiz{Knowledge not enough.} But intellectual acquirement and clearness are not enough. The good Catechist knows and mourns the fact that it is possible to be well informed and orthodox and valiant in truth, and yet to fail of
\quellseite{189}
% Continues the preceding page.
salvation. He realizes that to know about Christ is not yet life eternal. He wants this and something more. He wants the clear head always to be productive of the warm heart. Though I understand all mystery and all knowledge, and have not charity, or love, I am nothing. But love is a matter not of the head, but of the heart. Love is a feeling. Love is an experience. Love is life.

\randnotiz{Heart appropriation.} The good Catechist wants this experience, this life in his catechumens. He wants real heartfelt sorrow for and hatred of sin. He wants real heartfelt desire for forgiveness and for strength to overcome sin. He wants in each one a heartfelt assurance that his sins are forgiven; a heartfelt joy in this forgiveness, a heartfelt love to Him who has purchased and brought in this forgiveness. A heartfelt, personal appropriation of redemption is a prime requisite for Confirmation. Not that the experiences are to be all alike.\footnote{See the chapters on Conversion in “The Way of Salvation in the Lutheran Church.”} But all are to be real. None should be confirmed with a careless, cold or hostile heart. In so far as the Catechist is indifferent and negligent on this vital point, in so far he is not a good Lutheran Catechist. As has been noted before, and it cannot be emphasized too much, right here has been a sore spot in our catechetical practice. Our theory is beautiful and Scriptural.
\quellseite{190}
% Continues the preceding page.
Our practice has been sadly at fault in too many places. There is no such thing as a true Lutheran with a cold and loveless heart. Our Lutheran religion is essentially a heart-religion. To come to the Lord’s table with an impenitent and worldly heart is against the teaching and spirit of our Church. No catechumen, therefore, should be admitted to Confirmation and Communion without a living heart-religion. To bring it about the true Lutheran Catechist all through the course ceases not to labor, to travail, to wrestle, to pray; without some evidence of some measure of it; he would confirm no one.

\randnotiz{An obedient will.} The good Catechist knows that when he has the heart he has the will also. Where the heart has been moved and melted and won, the will says, “Lord, what wilt Thou have me to do?” The Catechist wants the will converted, turned, directed and ever ready to follow the Lord’s leading. Where the will is thus brought into subjection to the Lord’s will, there the Confirmation vow is a reality. Such confirmants can sing sweetly at the altar:
\begin{verse}
“My Jesus, as Thou wilt!\\
Oh, may Thy will be mine!\\
Into Thy hand of love\\
I would my all resign.\\
Through sorrow or through joy\\
Conduct me as Thine own,\\
And help me still to say,\\
‘My Lord, Thy will be done!’
\end{verse}

\quellseite{191}
% Continues the preceding page.
But how shall the Catechist know whether this desired fitness is present? How can he probe for it? How can he satisfy himself?

\randnotiz{Caution.} As we have sufficiently shown, we do not believe in an unnatural forcing of the emotions. Neither would we encourage a morbid introspection for frames and feelings. The good Catechist will do or say nothing that would do violence to the child-nature or encourage pretence, cant or hypocrisy. He will not be a Methodist. He will not fall into the unscriptural, fanatical, high-pressure methods of the revivalist. His sound theology will be in harmony with a sound psychology. Through the intellect he expects to reach the emotions and influence the will.

\randnotiz{Position of vantage.} And along this line he has been testing his class and each member of it all through the course. His teaching has been earnest and impressive. His applications and appeals have been direct and searching. Each pupil has been made to feel over and over again that the teacher is seeking his soul's good. The Catechist has noted the interest or lack of interest in each one. He has watched the desire and effort to get the lessons. He has seen where there has been attention, interest and earnestness. He has read the desires and the unspoken questions in the faces of his pupils. He has seen the whilom open and again concealed signs of emotion, inward
\quellseite{192}
% Continues the preceding page.
debate and life. He knows the general spiritual status of each one.

He has won the confidence and love of each one who can be won. He is looked up to and loved by all. He can, therefore, deal with each one now as he could not at the beginning of the course. The confidential relation has grown and ripened into this fair, filial fruitage.

\randnotiz{Personal interview.} The good Catechist is, therefore, in position to have a personal, private interview with each one. We strongly advise such an interview as the final test for Confirmation. But only that Catechist, who is a real consecrated \emph{seelsorger}, who deeply loves his Lord and these lambs of His purchase, can properly conduct this interview. The worldly pastor might do more harm than good. He would not know how to do it. He lacks the pastoral instinct. He does not know how to deal with the spiritual movements of the human soul. He would bungle at soul-curing.

At an appointed time, after the reviews discussed above, the good Catechist will have his class meet in a place where a private room is available. One by one he will take them into this private room. There he will speak kindly and tenderly to each one. He will invite perfect frankness and give the assurance that whatever is here disclosed will be sacredly kept and will never be repeated. The pastor who
\quellseite{193}
% Continues the preceding page.
would break this promise is not fit to be a pastor. He can ask the catechumen to tell why he desires Confirmation, what he knows and how he feels as to sin and sinning, as to sorrow for sin, trust in the Saviour and love to Him. He may ask as to habits of private prayer, feelings towards the Church, her doctrines and practices, and the private purpose towards abiding faithful.

All this must be done not in an inquisitorial or lordly manner, but gently, skillfully, paternally. The interview should be brief; where the class is large it must be very brief. In a very large class, it may be well to have two separate times for meeting and meet one-half of the class each time.

When the pastor has satisfied himself with one, let him quietly ask that one to kneel with him. There let him pour forth in calm, distinct and earnest words his own soul’s petitions for this \emph{beichtkind}, or spiritual child of his and commit him to the loving care and benediction of the Good Shepherd.

It is possible that even that child may go astray and become a prodigal. But he will never forget that interview and that prayer. The hope is that after such a scene, at the close of such a course of catechisation, as we have tried to outline and commend, it will be a rare exception for one to stray very far away. And if one does wander, in his
\quellseite{194}
% Continues the preceding page.
farthest wanderings these memories will haunt him, and there is hope that he will come to himself and come home again.

The writer of this can testify from personal experience and practice of this method, that it has proved itself to be fraught with untold blessing both for the pastor and for the child.

\quellseite{195}
\chapter[Chapter XVIII. The Public Examination]{Chapter XVIII. The Catechist and the Public Examination}

In many of our German and Scandinavian churches there is another service, preceding Confirmation, that is held in high esteem. So highly is it esteemed that these churches would not think of having a class confirmed without this service. Churches that do not have it are looked upon with suspicion. Their Lutheranism is seriously questioned. We refer to the public examination. 

\randnotiz{Public examination.} At an appointed time, either on the day of Confirmation or on a previous Sunday, the congregation is invited to come and witness the examination of the catechumens. These occupy the front seats and are supposed to be ready for the test on which their admission to Confirmation is to depend. Before the whole congregation, which on this day includes not merely the communicants, but all the friends and neighbors of the catechumens, the pastor now examines the class. He asks them collectively and individually to repeat the answers which they have committed to memory. He may throw in some Socratic questions to test their understanding of the catechetical material. But in the main the test has to do with that which has been memorized. This
\quellseite{196}
examination usually requires several hours. We know of one which lasted from ten in the morning till after three in the afternoon. 

\randnotiz{Historic custom.} This is an old historic custom. Traces of it are found in the early Church. In those days, when the catechumens were generally adult converts from Judaism and from heathenism, there was a special fitness in it. The congregations had a right to know and wanted to know whether the applicants for communicant membership understood and accepted the teachings and beliefs of the Church. They wanted both an intelligent acceptance and a pledge of fidelity from those who desired to be numbered with the saints. During the middle ages this testing was left to the priest, who was supposed to attend to it in private confession. 

The reformers and theologians of the Lutheran Church also favor the custom. Luther seems to take it for granted in his Preface to his Small Catechism. Chemnitz in his “Examen,” decidedly favors it.\footnote{Exam. Concil Trid., 1: 296.} The old Kirchenordnungen, or Church orders, give directions concerning it. The public examination is an old Lutheran custom. 

\randnotiz{Sound in principle.} Its grounds and principles are sound. It is in harmony with the Augsburg Confession: “For none are admitted, except they be first proved.”\footnote{Art. 24. See also Apol., Art. 24. Jacobs’ Book of Concord, Vol. I., pp. 50, 259.} With
\quellseite{197}
this principle we are in perfect harmony. We have shown above that we believe in the most rigid scrutiny before Confirmation. We have advised a double intellectual test, viz., one for the memory and the other for the understanding. We have advised that these several review examinations be so far public that the parents, guardians and special friends as well as the members of the Church Council be invited to be present. The parents and friends have a right to know what the children have learned during the course of instruction. The members of the Church Council are to advise with the pastor as to the admission of these candidates to communicant membership in the congregation. They ought to know something of the intellectual fitness of the applicants. In so far we believe in a public examination. 

We have also shown that we do not believe that the intellectual test is sufficient, something besides and in addition to intellectual equipment is necessary. This something cannot be publicly probed for or exhibited. This is the inner heart-fitness. For this the pastor alone must probe. And the result of this private probe will determine how many and who shall be admitted to Confirmation. Yes, we believe in examination, in more and deeper examination than do the advocates of the public test and exhibition before the congregation. We believe
\quellseite{198}
that our method of a twofold intellectual test, followed by a private probing, or, if you prefer, a real, frank, heart-to-heart private confession, is all sufficient. 

We have serious objections to the public examination as it is usually practiced. 

First. It easily becomes a temptation to pastor and pupil. \randnotiz{Temptation.} The pastor wants to make as good a showing as possible. A poor examination reflects on him. If the pupils do poorly the pastor gets at least a share of the blame. He may have done all he possibly could, and yet the Sunday display may be humiliating. This humiliation he desires to avoid. Therefore the temptation to cram for the occasion, to arrange for special answers from pre-appointed pupils, to make a show of attainment which is beyond the reality. The more private review and examination in the class-room is not open to these temptations. The public examination is likewise a temptation to the pupil. All sorts of quizzings and crammings and promptings are resorted to for the occasion. And if the pastor should stoop to a prearranged plan of assigning particular parts to particular persons, these become unconsciously dishonest—at the pastor’s suggestion. Such a temptation is not likely to come into the class-room review. 

Our second objection to the public examination is
\quellseite{199}
that it distracts the pupils from the main issue. \randnotiz{Distraction.} Their minds and hearts ought to be centered on the solemn step they are about to take. These closing days of the course ought to be filled with heart-searching and prayer and earnest meditation. But with the bugbear of the public examination before them the question that naturally becomes uppermost is, How will I stand the examination? Shall I pass? What if I fail? The very fear of possible failure, the nervous trepidation and excitement doubtless cause many undeserved failures. All this is absent from the pastor’s quiet and ordinary class-room review and test in the presence of a few invited and sympathizing friends. 

Our third, and even more serious, objection is that the public examination fosters the idea that \randnotiz{False standard.} intellectual fitness is the one thing needful. It encourages the idea that the brightest and readiest child is best prepared for Confirmation. We hesitate not to call such a principle orthodox rationalism. It has brought our beautiful and Scriptural custom of instruction and Confirmation into disrepute. It has wrought untold havoc in our dear Church. We do not claim or charge that all who favor and practice public examination harbor the pernicious principle that intellectual fitness is all-sufficient. But we do claim that this false principle is liable to creep in and has crept in, in thousands of places where the public
\quellseite{200}
examination is in vogue. Our method of classroom review and examination, followed by the private test or confession, obviates this serious danger. 

Our fourth and final objection is that even as an intellectual test the public examination is unfair. \randnotiz{Unfair test.} The timid, modest child may know the Catechism and may have come to answer well in class, and yet this faithful and promising pupil may become so embarrassed at the public examination as to utterly fail. On the other hand, the bold, forward pupil, who has not been nearly so faithful, may get far more credit than he deserves. This danger also is absent from the course we recommend. We believe, therefore, that all the good that may be claimed for the public examination can be conserved without it, and all the abuses to which it is liable to give rise can be avoided by our method. 

In the majority of our old English Lutheran congregations in the East the public examination has long since dropped out and is unknown. The Confirmation confession, pledge and vow are always made in public. The Confirmation service is impressive and solemn. It is complete without the public examination. 

We do not, however, advise our students or readers to be in undue haste to change old and established manners and customs. \randnotiz{Changing custom.} In a congregation where a public examination has been in vogue so
\quellseite{201}
long that it has become a part of the congregational life, it is better to make haste slowly. Let every effort be made to remove the false ideas and impressions that may be lodged among the people. Be very plain and persistent in showing both publicly and from house to house what does not and what does constitute fitness for Confirmation. Make it clear and drive it home that one may be well informed, sane and sound in intellect and yet be far from the kingdom of God. Rub it in over and over again that without heart-fitness, without experience, life and love, all else is worthless, sounding brass and tinkling cymbal. Thus make your people and your pupils see that the public examination does not settle the vital point, that it might be dispensed with without loss and with a possibility of much gain. Begin at once to pursue the method commended above. And thus, after you show them the more excellent way, they will not feel the need of the public examination.

\quellseite{202}
\chapter[Chapter XIX. Those Whom He Has Confirmed]{Chapter XIX. The Catechist Holding Those Whom He Has Confirmed}

In this last chapter on the science and art of catechising we come to the final serious problem. \randnotiz{Confirmed.} The regular meetings with the class are closed. The course of instruction and admonition is ended. The review examinations and the private interview have been held. The applicants who have been judged fit and worthy have been confirmed. The solemn and significant service has been held. The band of youth has encircled the altar. In earnest, trembling tones these young disciples have witnessed a good confession, pronounced their pledge and vowed allegiance to their Redeemer. Amid the prayers and tears of the many witnesses the pastor has laid his hands on each one and given to each one a personal benediction, than which he can give nothing more touching, tender and solemn. He has commended the young communicants to the prayerful interest, fellowship and helpfulness of the older members of the congregation. He has gone home from the service with conflicting thoughts and feelings. Happy is he if he can say, I have done my best, I have given my knowledge,
\quellseite{203}
my prayers, myself to the class. But, oh! how sad this Confirmation day and its subsequent reflections if he dare not seriously examine himself, lest he confess to himself that he has gone through the course as through an unwelcome task, that he has been all too often cold, careless, perfunctory, professional. 

But even when he has honestly and in the sight of God, with a heart burning with love for these young souls, done his prayerful best, what shall the end be with the newly confirmed? \randnotiz{Can all be held?} Will they be faithful unto death and in the end receive the crown of life? Can he hold them? Do pastors and churches generally hold the confirmed? What is the story of the Church records? To this question we referred in the Introduction to this book. It has come up again and again. It must be constantly before the conscientious pastor. Can he hold them? 

No, he cannot hold them all. Our Saviour did not hold all whom He had instructed. Not even all who had become professed disciples. “Many walked no more with Him.” Then said He to His disciples: “Will ye also go away?” The Apostles and inspired evangelists did not hold all. The best Catechists have never held all. The most gifted and consecrated preachers in our own and in other churches have never held all whom their preaching awakened and won.

\quellseite{204}
\randnotiz{Why not?} Why do so many lapse? Doubtless chiefly because of the sinfulness and treacherousness of the human heart. Doubtless also because of the faults in the Catechists and the catechising referred to and warned against in almost every chapter of this book. With the lessons of this book taken to heart and put into practice, the losses will be greatly reduced. This, at least, is the purpose, the prayer and the hope that inspired the writing of the book. But this is a digression. 

Many of the losses are also occasioned by the neglect or the wrong treatment after Confirmation. More can be held if the proper treatment is now given. What, then, can the good pastor do to hold those whom he confirms? A few helpful hints are submitted. 

First. Before and after Confirmation it must be made clear and impressed that the congregation has a duty and a responsibility towards its new communicants. \randnotiz{Members of congregation can aid.} Collectively and as individuals the older members must realize their relation and their duty to the incoming ones. The congregation is a family. It is constantly adding to and adopting children into the household of faith. The family of the redeemed should naturally show a warm hospitality toward the newcomers. Many of them are strangers, others are timid, still others weak and wavering. They are more or less dazed by their new relation. They do not know how to take it.

\quellseite{205}
Now, let the older members give the warm hand of Christian fellowship and the kind word of welcome. Let a public, informal reception be held for them. At this reception let the older members welcome the new, introduce them to others, invite them to their homes and so make them realize that they have come into a goodly company, that the lines have fallen to them in pleasant places, that theirs is a blessed heritage. And let the pleasant greetings and kind words be kept up every Sunday. The Apostolic Church had a warm-hearted and warm-handed social side. Our churches are too often cold, exclusive, repellent towards strangers. It might also be suggested that the boy is in special need of attention and companionship. The boy likes to be noticed. The boy is too often overlooked and ignored by church members. The boy feels it. The members of the congregation can do much towards holding the newly confirmed. But here also the pastor must instruct, show and lead the congregation. 

Second. Another help towards holding the young is to organize them into an association. We know the danger of over-organizing the church. We deprecate the splitting up and dividing until, instead of one congregation, the pastor has on his hands a multitudinous array of separate and separatistic organizations in an unorganized aggregation of
\quellseite{206}
people. \randnotiz{Organize association.} And yet our age demands some organization within the one organization. The danger also can be minimized, if not entirely averted, if the pastor makes each organization with all its members realize that it is a department of the congregation, or, if you will, a large committee of the church, with a specific part of the church’s work assigned to it. The solidarity of the congregation must ever be kept uppermost. With these precautions there can be an association or union of the confirmed, divided into classes, somewhat analogous to the alumni associations of our schools. Let the association have a very simple constitution, by-laws and order of exercises. Let it be properly officered and have the necessary committees. Among them let there be an outlook committee, which will keep tab on the church attendance and activity of all the members, reporting the absentees and the endangered ones. Another committee should see to it that all such are gone after and reclaimed if possible. They should also be reported to the pastor, who should gladly assist in lining up delinquents and reclaiming wanderers. Meetings should be held at least twice a year. At these meetings the officers and committees should read carefully prepared reports, the roll should be called, all absentees and careless ones should be reported, inquired after and counsel taken as to their reclamation. The pastor should keep in
\quellseite{207}
close touch with the organization, encourage and help them to keep up a laudable esprit de corps for their congregation and church. No doubt such an organization could do much towards holding the young. In large congregations each class might have a sub-organization to work under and in harmony with the general body. 

\randnotiz{Enroll in Sunday school.} Third. The pastor should see to it that every confirmed one should be enrolled in his proper class in the graded Sunday school. The Sunday school teacher should also take a special interest in these members of the class and do everything possible to keep them interested and faithful. 

\randnotiz{In Luther League.} Each one should also be enrolled as a member of the Luther League. As it is one of the constitutional objects of this useful organization to keep its members loyal to their Confirmation vows, it goes without saying that every League will do all it can to interest and hold such members. Thus the Sunday school and League can work together to help the pastor to hold those whom he has confirmed. 

\randnotiz{Adapted service.} Fourth. Let the pastor, the association of the confirmed, the Sunday school and the Luther League all join in trying to find and give work to the careless and unsafe ones. Make each such an one realize that he is “part of the concern.” Throw on him some definite, specific responsibility, that will require his presence and effort. But be careful not
\quellseite{208}
to try to compel him to do that which is obnoxious to him. Some are driven out of the organizations by constant urging to write and read papers, open and lead meetings, and do such things for which they have a positive aversion. Find something that is easier and more congenial for them, and select those who have an aptitude for these harder tasks. But find or invent some kind of work for the laggards. This applies also to their place in the congregation. As soon as they are old or big enough find places for them in the choir, or as ushers, librarians, assistants of one kind or another. Many of them will be found to appreciate office and responsibility. And many an one can thus be held, who would otherwise stray away. 

Fifth. Along with all this, in addition to it all and above all, the pastor must keep up an active, untiring personal and pastoral interest in each one whom he has confirmed. \randnotiz{Pastor’s personal interest.} He ought to know, by this time, the weaknesses, the temptations, the dangers of each one. Careful and constant pastoral visiting is needed. Each one must constantly realize that the pastor is his best friend and is ever seeking his highest good. The pastor can often see the dangers and the needs better than the parents. An active and energetic boy may be in danger of ruin through idleness. For his physical, moral and spiritual salvation that boy needs work—regular daily work.

\quellseite{209}
The pastor needs to impress this upon the parents and should be ready to lend his kindly offices in helping to find suitable employment. In other cases poverty may make it necessary for the boy or girl to find employment. And here again the good pastor will be glad to help. And all the while the good pastor will keep up what he began in the catechetical class, viz., the directing of the gifted ones and the winning of them for the higher vocations, the ministry of mercy and the ministry of the Word. He will ever strive with the youth and with their parents that they may be given an opportunity for all the higher Christian education they can get in the institutions of their own church. Much catechising and counseling of the parents will be required in order that they may not instill into their children the baleful doctrine that money is the chief thing in life. These parents need to be converted to the idea that there are some things that are better than money, that Christian education, Christian character and Christian service are of more worth than the gold of the Indies without them. Many a bright and promising youth has been weaned away from the Church and her Christ by the glitter of gold, the love of pleasure and the glamor of fame. True \emph{seelsorge} will save many an one from such snares of sin and soul suicide.

The faithful pastor will make every effort to keep
\quellseite{210}
\randnotiz{Care for those who remove.} trace of every confirmed one who removes from the bounds of his parish. He will, if possible, counsel such an one before he goes as to the dangers and safeguards in his new home. He will follow the departed ones with kindly letters; will, if possible, introduce them to helpful friends, direct them to the nearest church of their own faith, write the pastor of that church and have him look up and care for the stranger, and so help to hold the absent one in another true fold when he cannot keep him in his own. All this is pastoral theology. It is private \emph{seelsorge}. It is of vital importance in holding the confirmed catechumens. Its practice would largely lessen our sad losses. It would remove one of the frequent objections to our method of gathering the young into the communion of the Church. It would make many a congregation renew her youth like the eagle. It would bring the beauty of the Lord our God upon us. He would establish the work of our hands upon us; yea, the work of our hands, He would establish it.

Yes, we verily believe that if we could have the right kind of catechising, confirming and care after Confirmation the time would speedily come when our dear Church would go forth as the morning, bright as the sun, fair as the moon and terrible as an army with banners.

\emph{To this end help us, good Lord!}
